\documentclass[a4paper]{article}
\usepackage{amsfonts}
\usepackage{amsmath}
\usepackage{amssymb}
\usepackage{graphicx}

\begin{document}
  
\noindent \large Auteur: Michael Livchits \\
\noindent \large Studierichting: Informatica\\
\noindent \large Oefeningenreeks 7A nr. 16.4\\

\medskip

\normalsize

$f(x,y) = 
\left\{ 
\begin{array}{ll} 
\dfrac{x^2y^8}{x^6+y^{12}} & \text{als } (x,y) \neq (0,0)\\[2ex] 
0 & \text{als } (x,y) = (0,0) 
\end{array} 
\right.$ \\

Continuiteit:\\

Stel $x = y^2$\\

$\lim\limits_{y\rightarrow 0}f(y^2,y) = \dfrac{y^{12}}{y^{12}+y^{12}}$\\

$= \dfrac{1}{2}$\\

Nu voor $x = 0$\\

$\lim\limits_{y\rightarrow 0}f(0,y) = \dfrac{0y^{8}}{0+y^{12}}$\\

$= 0$\\

$\Rightarrow f$ is niet continu in $(0,0)$\\

Partiele afgeleiden:\\

$\dfrac{\partial f}{\partial x}(0,0) = \lim\limits_{\lambda \rightarrow 0} \dfrac{f(\lambda, 0) - f(0,0)}{\lambda} = \lim\limits_{\lambda \rightarrow 0}\dfrac{0}{\lambda^6} = 0$\\

$\dfrac{\partial f}{\partial y}(0,0) = \lim\limits_{\lambda \rightarrow 0} \dfrac{f(0, \lambda) - f(0,0)}{\lambda} = \lim\limits_{\lambda \rightarrow 0}\dfrac{0}{\lambda^{12}} = 0$\\

$\Rightarrow f$ is partieel afleidbaar in $(0,0)$\\

Afleidbaarheid:\\

Stel $\overline{h} = (h_1,h_2)$, $h_1\neq0\neq h_2$\\

$Df(\overline{0}, \overline{h}) = \lim\limits_{\lambda \rightarrow 0} \dfrac{f(\lambda h_1, \lambda h_2) - f(0,0)}{\lambda}$\\

$= \lim\limits_{\lambda \rightarrow 0} \dfrac{\dfrac{\lambda^{10}h_1^2h_2^8}{\lambda^6h_1^6+\lambda^{12}h_2^{12}}}{\lambda}$\\

$= \lim\limits_{\lambda \rightarrow 0} \dfrac{\lambda^{10}h_1^2h_2^8}{\lambda^7(h_1^6+\lambda^6h_2^{12})}$\\

$= \lim\limits_{\lambda \rightarrow 0} \dfrac{\lambda^3h_1^2h_2^8}{h_1^6+\lambda^6h_2^{12}}$\\

$= 0$\\

$\Rightarrow f$ is afleidbaar in $(0,0)$\\

Differentieerbaarheid:\\

$f$ is niet continu in $(0,0)$\\

$\Rightarrow f$ is niet differentieerbaar in $(0,0)$\\

\begin{thebibliography}{999}
%\bibitem{a} Referentie
\end{thebibliography}
\end{document} 